
\subsubsection{4.1 Research Questions}

Three research questions organize the experimental analysis.

\textbf{RQ1}: Do AIFS features predict annotator preference in a conditional logit model that controls for response length and prompt-level heterogeneity via cluster fixed effects?

\textbf{RQ2}: Do annotators in the online (deployed-AI) condition show a stronger conditional preference for AIFS features relative to annotators in the base condition, as measured by the interaction coefficient β₄?

\textbf{RQ3}: If the result for RQ2 is null, can insufficient statistical power be ruled out as the explanation, and what alternative accounts are consistent with the data?

RQ1 serves as a construct validity check: a significant β₁ establishes that AIFS captures preference-relevant variation, which is a precondition for interpreting any null result on β₄ as substantively informative. RQ3 is addressed through confidence interval analysis and sensitivity testing.

\subsubsection{4.2 Dataset Summary}

\begin{table}[htbp]
\centering
\begin{tabular}{ll}
\toprule
Property & Value \\
\midrule
Dataset & Anthropic HH-RLHF [Bai et al., 2022] \\
Helpful-base rows (total / after filter) & 43,835 / 20,108 \\
Helpful-online rows (total / after filter) & 22,007 / 20,063 \\
Total preference pairs & 80,342 \\
Embedding model & all-MiniLM-L6-v2 [Reimers \& Gurevych, 2019] \\
Clustering threshold (cosine similarity) & 0.85 \\
Unique semantic clusters & 27,034 \\
Median cluster size & 2.8 pairs \\
Mean cluster size & 3.0 pairs \\
\bottomrule
\end{tabular}
\end{table}

\subsubsection{4.3 AIFS Feature Extraction}

For each response, AIFS is computed as the sum of four binary regex pattern indicators: (1) structured_list, (2) safety_preface, (3) cot_marker, (4) hedging. The within-pair AIFS contrast (ΔAIFS = AIFS_chosen − AIFS_rejected) is the primary treatment variable. Computing the contrast at the pair level removes prompt-level AIFS demand, isolating supply-side formatting variation.

\subsubsection{4.4 Model Specification}

```
P(chosen = 1) ~ β₁·ΔAIFS + β₂·Δlength + β₄·(ΔAIFS × split) + cluster_FE
```

All cluster fixed effects are estimated by conditional maximum likelihood. The optimizer is Newton-Raphson (BFGS failed to converge on the full dataset). The implementation uses the \texttt{statsmodels} ConditionalLogit class with \texttt{groups=cluster_id}.

\subsubsection{4.5 Evaluation Metrics}

\begin{table}[htbp]
\centering
\begin{tabular}{lll}
\toprule
Metric & Role & Pre-registered Threshold \\
\midrule
β₄ point estimate & Directional check & > 0 \\
OR = exp(β₄) & Effect magnitude & ≥ 1.10 \\
95\% CI lower bound & Precision / exclusion of null & > 1.0 \\
Wald p-value & Statistical significance & < 0.01 \\
LRT statistic & Model fit contribution & Complementary \\
\bottomrule
\end{tabular}
\end{table}

